% !PRINTED 329
\weberpage{329}
%<<CONT
time by rows, the other time by columns, it follows by the
formulas (4):

\weq{5}{C_1\,(y,b) = r^6\,C_1\,(x,a),}

whereby the invariant property is expressed and the weight
determined to be $= 6$.

The third covariant $C_2$ that we consider is the functional determinant
 of the three forms $f$, $H$, $C_1$, thus

\weq{6}{C_2\,(x, a) = \begin{vmatrix} f'\,(x_1), & H'\,(x_1), & C_1'\,(x_1) \\ f'\,(x_2), & H'\,(x_2), & C_1'\,(x_2) \\ f'\,(x_3), & H'\,(x_3), & C_1'\,(x_3) \end{vmatrix},}

which, as already follows from \S.\,59 of the first volume, has the
invariant property:

\weq{7}{C_2\,(y,b) = r^9\,C_2\,(x,a),}

and which therefore has the weight 9.

If we denote by $\lambda$ the weight, by $\nu$ the degree in the
variables, by $\mu$ the degree in the coefficients, the following
compilation results:

\begin{center}
\begin{tabular}{lccc}
 & $\nu$, & $\mu$, & $\lambda$ \\
\hline
$f$ & $n$, & 1, & 0 \\
$H$ & $3n-6$, & 3, & 2 \\
$C_1$ & $8n-18$, & 8, & 6 \\
$C_2$ & $12n-27$, & 12, & 9.
\end{tabular}
\end{center}

\begin{center}
\S.\,88.

\gesperrt{The Hesse Curve.}
\end{center}

The first of the covariants introduced, $H$, has a simple
geometric meaning, with which we now wish to become acquainted.
This investigation is essentially facilitated by the fact that,
on account of the invariant property of the form $H$, we lose
nothing in generality if we give the coordinate system any
special position with respect to the curve $f$.

If we place the vertex $\xi_3$ of the coordinate triangle at a
point of the curve $f$ that does not belong to the singular ones, and
choose the tangent at this point as the $x_1$-axis, then the form $f$
takes, by \S.\,86, (3), the expression:

\weq{1}{f = h\,x_3^{n-1}x_1 + x_3^{n-2}(a\,x_1^2+2\,b\,x_1x_2+c\,x_2^2)+x_3^{n-3}f_3+\cdots,}
%CONT>>
